%=====================================================================
\section{The demand side}\label{sec:dem:top}
\label{sec:dem:interface}

Everything upstream of this section --- the well, the water circuit, the coil, the
wheel, the cooler --- exists to deliver one air stream to one space. This section
defines the only thing the plant is allowed to know about that space.

The model's demand side is deliberately \emph{generic}. The building is not a
hatchery, a data hall or an office as far as the plant is concerned: it is a pair
of numbers per time step, a sensible load and a moisture load, plus one closure
that says how much air is available to carry them. A hatchery is one way of
producing that pair; Section~\ref{sec:dem:hatch} works it out in full, but it is
an \emph{instance}, not the definition. If a second application ever replaces it,
nothing in Sections~1--5 of this report changes.

\subsection{Two balances, three unknowns, one closure}
\label{sec:dem:balances}

\paragraph{Hypotheses, stated before the algebra.} The space is treated as a
single steady, well-mixed control volume at constant pressure $\patm$, served by one
supply stream of dry-air mass flow $\md_{\mathrm{sup}}$ and returning at the room
condition. ``Well mixed'' means the extract state \emph{is} the room state
$(T_{\mathrm{pc}},w_{\mathrm{pc}})$: no stratification, no local gradients, no
distribution effectiveness. ``Steady'' means the room's thermal and moisture
capacitance is ignored, so the balance has no time derivative and each time step
is independent of the last.

\paragraph{Dry-air conservation.} With no infiltration \emph{through the supply
path}, the dry-air flow leaving equals the flow entering, so a single
$\md_{\mathrm{sup}}$ appears in both balances below. This is why humidity ratio
(per kilogram of \emph{dry} air) is the right moisture variable: it makes the
water balance linear in a conserved flow.

\paragraph{Water balance.} Let $\mdw$ be the net rate at which water vapour is
released into the space by everything that is not the supply air --- occupants,
process, infiltration, evaporating product. Steady state requires that the supply
air carry it away:
\begin{equation}
\md_{\mathrm{sup}}\,w_{\mathrm{pi}}+\mdw=\md_{\mathrm{sup}}\,w_{\mathrm{pc}}
\qquad\Longleftrightarrow\qquad
\boxed{\;\mdw=\md_{\mathrm{sup}}\left(w_{\mathrm{pc}}-w_{\mathrm{pi}}\right)\;}
\label{eq:dem:water}
\end{equation}

\paragraph{Energy balance.} Likewise, with $\qs$ the net sensible gain of the
space and $c_p$ the moist-air specific heat per kilogram of dry air,
\begin{equation}
\md_{\mathrm{sup}}\,c_p\,T_{\mathrm{pi}}+\qs=\md_{\mathrm{sup}}\,c_p\,T_{\mathrm{pc}}
\qquad\Longleftrightarrow\qquad
\boxed{\;\qs=\md_{\mathrm{sup}}\,c_p\left(T_{\mathrm{pc}}-T_{\mathrm{pi}}\right)\;}
\label{eq:dem:sens}
\end{equation}
The code evaluates $c_p$ at the \emph{room} humidity by default,
$c_p=(1.006+1.86\,w_{\mathrm{pc}})\times10^{3}\;\si{\joule\per\kilogram\per\kelvin}$,
and lets the caller override it. Evaluating at the room rather than the supply
state is a choice of linearisation point; at the design case of
Section~\ref{sec:dem:worked} the two differ by less than a part in $10^{3}$.

\paragraph{Sign convention --- fix it once and never revisit it.}
Both loads are \emph{positive when the plant must remove something}:
\begin{center}
\begin{tabular}{@{}lll@{}}
\toprule
Quantity & Positive means & Consequence for the supply state \\
\midrule
$\qs>0$   & cooling required        & $T_{\mathrm{pi}}<T_{\mathrm{pc}}$: supply colder than the room \\
$\mdw>0$  & dehumidification required & $w_{\mathrm{pi}}<w_{\mathrm{pc}}$: supply drier than the room \\
\bottomrule
\end{tabular}
\end{center}
Negative values are perfectly legal and mean heating and humidification: the
same two equations then return a supply state \emph{above} the room in
temperature or humidity. Nothing in the interface branches on the sign, which is
what makes the plant capable of running in all four quadrants without a mode
switch.

\paragraph{The counting problem.} Equations~\eqref{eq:dem:water} and
\eqref{eq:dem:sens} are two equations in \emph{three} unknowns:
$\md_{\mathrm{sup}}$, $T_{\mathrm{pi}}$ and $w_{\mathrm{pi}}$. The loads are
data; the supply is not determined. Exactly one further statement --- a
\emph{closure} --- is required, no more and no less. The code enforces this
literally: supplying both closures, or neither, raises an error rather than
silently preferring one.

\begin{center}
\fbox{\parbox{0.92\textwidth}{
\textbf{The two admissible closures.}\\[2pt]
\emph{Constant volume.} Prescribe the flow $\md_{\mathrm{sup}}$. The supply
temperature and humidity then follow from
Eqs.~\eqref{eq:dem:water}--\eqref{eq:dem:sens} and float with the load. This is
the fixed-fan machine: the air handler moves a fixed mass and the conditioning
does all the work.\\[3pt]
\emph{Variable volume.} Prescribe the supply depression
$\Delta T_{\mathrm{sup}}=T_{\mathrm{pc}}-T_{\mathrm{pi}}$. The flow then follows
the sensible load,
$\md_{\mathrm{sup}}=\max\!\left(|\qs|/(c_p|\Delta T_{\mathrm{sup}}|),\,
\md_{\min}\right)$, and the supply temperature is held. This is the VAV machine:
the coil holds a discharge setpoint and the fan tracks the load.
}}
\end{center}

Two details of the variable-volume branch are worth reading off the code. First,
the absolute values: the depression is treated as a \emph{magnitude}, so a
heating load ($\qs<0$) returns a positive flow and, through
Eq.~\eqref{eq:dem:sens}, a supply \emph{above} room temperature. Second, the
floor $\md_{\min}$: a minimum flow (typically a mandatory fresh-air rate) can be
imposed, and when it binds the interface returns a \verb|vent_limited| flag
rather than quietly delivering less air than the application requires.

\paragraph{The room condition line.} Divide Eq.~\eqref{eq:dem:water} by
Eq.~\eqref{eq:dem:sens}. The flow --- the one thing the closure supplies ---
cancels identically:
\begin{equation}
\frac{w_{\mathrm{pc}}-w_{\mathrm{pi}}}{T_{\mathrm{pc}}-T_{\mathrm{pi}}}
=\frac{c_p\,\mdw}{\qs}\;\equiv\;S .
\label{eq:dem:rcl}
\end{equation}
This is the classical \emph{room condition line}: on a psychrometric chart, the
locus of admissible supply states is a straight line through the room point
$(T_{\mathrm{pc}},w_{\mathrm{pc}})$ whose slope $S$ is set by the load
\emph{ratio} alone.

\begin{center}
\fbox{\parbox{0.92\textwidth}{
\textbf{What is fixed by what.} The load ratio fixes the \emph{direction} in
which the supply state must lie from the room point. The closure fixes \emph{how
far along that line} it sits --- a large flow puts the supply near the room, a
small flow pushes it far away. Direction is physics; distance is a design
choice. Every disagreement about ``the required supply air'' is one of these two
and it is worth naming which.}}
\end{center}

\paragraph{Why the sensible heat ratio cannot classify the load.} It is
tempting to compress the load pair into the single number reported by the
interface,
\begin{equation}
\mathrm{SHR}=\frac{\qs}{\qs+\ql},\qquad \ql=\mdw\,\hfg ,
\label{eq:dem:shr}
\end{equation}
and to read the operating quadrant off it. \textbf{This does not work, and the
failure is not subtle.} The ratio is invariant under a simultaneous sign flip of
both loads: $\qs=+100$, $\ql=+25$ and $\qs=-100$, $\ql=-25$ both give
$\mathrm{SHR}=0.8$. Two \emph{negative} loads --- a space that needs heating and
humidification, the deep-winter quadrant --- return a perfectly respectable
positive SHR indistinguishable from summer cooling. A ratio destroys exactly the
information (the signs) that selects the quadrant. The interface therefore
classifies on the signs of $\qs$ and $\mdw$ \emph{separately} and reports
SHR as a diagnostic only.

\paragraph{What the steady, well-mixed hypothesis costs.} Three things. There is
no thermal mass, so a load spike is passed instantly to the plant instead of
being buffered by the structure --- the model cannot exploit pre-cooling or
night flush, and it over-states peak duty on a short excursion. There is no
stratification, so a tall space that could be conditioned only in the occupied
layer is charged for its whole volume. And there is no room control dynamics:
the space is \emph{assumed} to be held at
$(T_{\mathrm{pc}},\RH_{\mathrm{pc}})$, so the model answers ``what must be
supplied'' and never ``what would the room drift to if the plant fell short''.
The saturation flags elsewhere in the model exist precisely because that second
question has been assumed away here.

\subsection{Why the latent load is a moisture rate and not a power}
\label{sec:dem:hfg}

The interface carries the latent load as $\mdw$ in \si{\kilogram\per\second}, not
as a power in \si{\kilo\watt}. This looks like a units preference. It is not; it
removes a silent error.

\paragraph{The double conversion.} A load generator that measures evaporation ---
which is what physically happens --- knows a \emph{rate}. To report kilowatts it
must multiply by some latent heat, $\ql=\mdw\hfg^{(1)}$. The interface, needing a
rate for Eq.~\eqref{eq:dem:water}, must then divide by some latent heat,
$\mdw^{\ast}=\ql/\hfg^{(2)}$. The physical quantity therefore passes through
$\hfg$ \emph{twice}, and only if the two values agree exactly does the round trip
return what was sent:
\begin{equation}
\mdw^{\ast}=\mdw\,\frac{\hfg^{(1)}}{\hfg^{(2)}} .
\label{eq:dem:hfground}
\end{equation}

\paragraph{The error is not hypothetical.} Reasonable values of $\hfg$ near room
temperature span roughly \SIrange{2430}{2500}{\kilo\joule\per\kilogram} depending
on whether one quotes it at \SI{25}{\celsius}, at \SI{0}{\celsius}, or as the
enthalpy datum embedded in the moist-air enthalpy correlation. The present
notebook itself carries two: the load reference
$\hfg=\SI{2.5}{\mega\joule\per\kilogram}$ used by the interface and every load
generator, and the datum \SI{2501}{\kilo\joule\per\kilogram} inside the moist-air
enthalpy function (and inside the legacy \verb|building_loads| helper). They
differ by \SI{1}{\kilo\joule\per\kilogram} here, harmlessly --- but nothing in a
kilowatt-valued interface would have detected it had they differed by fifty.

\paragraph{Size of the damage.} Propagating Eq.~\eqref{eq:dem:hfground} through
Eq.~\eqref{eq:dem:water}, the induced error falls entirely on the required supply
\emph{depression}:
\begin{equation}
\frac{\delta\!\left(w_{\mathrm{pc}}-w_{\mathrm{pi}}\right)}
     {w_{\mathrm{pc}}-w_{\mathrm{pi}}}
=\frac{\delta \hfg}{\hfg}
\;\;\approx\;\;\frac{50}{2500}=2\,\% .
\label{eq:dem:hfgerr}
\end{equation}
Two per cent of a humidity difference is far below the resolution of any humidity
sensor and far below the scatter of the wheel model that must deliver it. It will
never be caught by inspection of a result, and it will never be caught by a unit
check either, because kilowatts are kilowatts. \textbf{An error that is both
systematic and invisible is worse than a large one.} The cure is structural, not
numerical: convert once, at the boundary, and keep the conserved quantity ---
mass --- in the interior.

\paragraph{How the code implements it.} A single module-level constant
\verb|LOAD_HFG| is the only conversion the demand side may use. The generic
generator \verb|prescribed_loads| exists precisely to serve callers who genuinely
have kilowatts: it accepts $\ql$ in \si{\kilo\watt} together with an
\emph{optional explicit} $\hfg$, converts there and then, and hands a rate
onwards. The kilowatt figure the interface reports back, $\ql=\mdw\,$\verb|LOAD_HFG|,
is a derived display quantity and is never used to compute anything.

\subsection{The guards}
\label{sec:dem:guards}

The two balances are linear and always return \emph{a} solution. The question is
whether that solution is a state that air can be in. Two failures used to pass
silently and are now trapped.

\paragraph{Guard 1: negative required supply humidity.} Equation~\eqref{eq:dem:water}
gives $w_{\mathrm{pi}}=w_{\mathrm{pc}}-\mdw/\md_{\mathrm{sup}}$, which goes
negative when
\begin{equation}
\md_{\mathrm{sup}}<\frac{\mdw}{w_{\mathrm{pc}}} .
\label{eq:dem:guard1}
\end{equation}
Physically: \emph{even bone-dry supply air, at that flow, cannot carry the
moisture away}. The space would have to be served air with negative water
content. Equation~\eqref{eq:dem:guard1} is therefore an absolute lower bound on
the supply flow imposed by the latent load alone, independent of any equipment,
and it is worth evaluating before sizing anything: no desiccant, however good,
buys a way around it.

\paragraph{Guard 2: supersaturated required supply state.} Even a positive
$w_{\mathrm{pi}}$ can be unreachable. The supply temperature is fixed
independently by Eq.~\eqref{eq:dem:sens}, so the pair
$(T_{\mathrm{pi}},w_{\mathrm{pi}})$ may fall \emph{below the saturation line}.
The interface inverts the humidity ratio to a vapour pressure and tests it
against the saturation pressure at the required supply temperature:
\begin{equation}
\RH_{\mathrm{pi}}
=\frac{1}{\psat(T_{\mathrm{pi}})}\cdot
\frac{w_{\mathrm{pi}}\,\patm}{0.622+w_{\mathrm{pi}}}\;>\;1
\quad\Longrightarrow\quad\text{no component can deliver it.}
\label{eq:dem:guard2}
\end{equation}
This is not an equipment shortfall, it is an inconsistency between the load pair
and the closure: the requested state does not exist. The usual cause is a small
supply flow with a large sensible load and a modest latent load --- the supply is
driven very cold while still being asked to stay humid. The remedies are to raise
the flow (moving the supply state back toward the room along the line of
Eq.~\eqref{eq:dem:rcl}) or to accept reheat, which means deliberately leaving the
room condition line.

\paragraph{Reading the two guards together.} Geometrically they fence the room
condition line at both ends. The admissible supply state must lie on the line of
Eq.~\eqref{eq:dem:rcl}, above $w=0$ and below the saturation curve
$w=w_{\mathrm{sat}}(T)$. The closure slides the state along that line; the guards
say when it has slid off the end of the physically available segment.

\paragraph{A deliberate asymmetry in severity.} Malformed \emph{inputs} raise
exceptions and stop the run: both closures given or neither, a zero depression, a
non-positive flow. Physically impossible \emph{outputs} raise warnings and let
the run continue. That is a considered choice: a design sweep that crashes on its
first infeasible corner tells you nothing about the shape of the feasible region,
whereas one that flags and continues maps it. The cost is that the flags must
actually be read --- a warning nobody looks at is worse than an exception.

%=====================================================================
\section{A worked demand case: the poultry hatchery}
\label{sec:dem:hatch}

The hatchery is now one implementation of the generator contract of
Section~\ref{sec:dem:interface}: a function of ambient state returning the pair
$(\qs,\mdw)$. It is retained because it is the project's motivating application
and because it exercises the interface hard --- an internal heat source, an
internal moisture source, a room \emph{warmer} than the Gulf ambient, and a
mandatory ventilation rate that the occupants (embryos) set and the designer
cannot negotiate.

\subsection{Internal terms}
\label{sec:dem:internal}

\paragraph{Embryonic metabolic heat.} Developing embryos respire and the heat
release rises steeply toward hatch. Written per egg,
\begin{equation}
Q_{\mathrm{eggs}}=N_{\mathrm{eggs}}\,q_{\mathrm{egg}},
\qquad q_{\mathrm{egg}}=\SI{0.13}{\watt}\ \text{per egg},
\label{eq:dem:Qeggs}
\end{equation}
a \emph{sensible source}: it flows into the space. Note how this inverts the
usual comfort-cooling intuition, in which the envelope is the load and the
interior is passive.

\paragraph{Egg moisture loss.} A setter deliberately drives controlled moisture
loss from the eggs --- of order \SI{12}{\percent} of egg weight over incubation,
and essential to hatchability. That water evaporates into the room air:
\begin{equation}
\dot m_{\mathrm{evap}}=N_{\mathrm{eggs}}\,r_{\mathrm{evap}},
\qquad r_{\mathrm{evap}}=\SI{4.6e-9}{\kilogram\per\second}\ \text{per egg}
\;(\approx\SI{0.40}{\gram}\ \text{per egg per day}).
\label{eq:dem:mevap}
\end{equation}
The application therefore \emph{wants} the eggs to lose water, which is why the
room humidity cannot be allowed to drift up: too humid means insufficient weight
loss and a poor hatch. The humidity setpoint here is a product specification, not
a comfort preference.

\subsection{Envelope terms in conductance form}
\label{sec:dem:envelope}

\paragraph{Sensible: two genuine conductances.} Both external sensible paths are
linear in a temperature difference, so both can be written as a conductance times
a driving potential:
\begin{equation}
Q_{\mathrm{sens,ext}}
=\underbrace{H_{\mathrm{tr}}}_{\text{material}}\left(T_{\mathrm{sol}}-T_{\mathrm{pc}}\right)
+\underbrace{\dot m_{\mathrm{inf}}\,c_{p,\mathrm{a}}}_{\text{air mass}}
\left(T_{\mathrm{oa}}-T_{\mathrm{pc}}\right),
\qquad H_{\mathrm{tr}}=\sum_i U_iA_i .
\label{eq:dem:Qsensext}
\end{equation}
$H_{\mathrm{tr}}$ [\si{\watt\per\kelvin}] is a property of materials and
geometry; $\dot m_{\mathrm{inf}}c_{p,\mathrm{a}}$ [\si{\watt\per\kelvin}] is a
property of the airtightness. The infiltration here is \emph{uncontrolled
leakage} and is distinct from the deliberate ventilation, which is handled in the
process loop (Section~\ref{sec:dem:assembled}).

\paragraph{Sol-air temperature.} For the opaque envelope the driving outdoor
temperature is not the air temperature: absorbed solar radiation raises the
surface above it. The standard device folds this into an equivalent outdoor
temperature,
\begin{equation}
T_{\mathrm{sol}}=T_{\mathrm{oa}}
+\frac{\alpha_s I_{\mathrm{sol}}}{h_o}-\frac{\varepsilon\,\Delta R}{h_o}
\;\;\longrightarrow\;\;
T_{\mathrm{sol}}=T_{\mathrm{oa}}+\Delta T_{\mathrm{solar}},
\qquad \Delta T_{\mathrm{solar}}=\SI{4}{\kelvin},
\label{eq:dem:solair}
\end{equation}
with $\alpha_s$ the solar absorptance, $I_{\mathrm{sol}}$ the incident
irradiance, $h_o$ the outside film coefficient and the last term the long-wave
sky correction. The model currently uses the constant offset on the right; under
hourly weather forcing it becomes the full expression on the left. For a
single-storey hatchery, whose roof is most of its envelope, this term is not
cosmetic.

\subsection{The latent asymmetry, which is the physically central point}
\label{sec:dem:asymmetry}

The external latent load has \emph{one} term and no counterpart to
$H_{\mathrm{tr}}$:
\begin{equation}
\dot m_{w,\mathrm{ext}}=\dot m_{\mathrm{inf}}\left(w_{\mathrm{oa}}-w_{\mathrm{pc}}\right).
\label{eq:dem:mwext}
\end{equation}

\begin{center}
\fbox{\parbox{0.92\textwidth}{
\textbf{The asymmetry, stated plainly.} A wall conducts heat but not vapour.
There is no material vapour conductance analogous to a $U$-value: vapour
diffusion through permeable construction exists but is second order and is
neglected here. Hence the \emph{sensible} external load has two knobs ---
insulation ($H_{\mathrm{tr}}$) and tightness ($\dot m_{\mathrm{inf}}$) --- while
the \emph{latent} external load has only the second.\\[3pt]
\textbf{In a humid climate, insulation does nothing for humidity.} Airflow
control is the entire latent game. When a reference speaks of a ``latent
conductance'' it means the air-mass term $\dot m\,\hfg$, never a material
property.}}
\end{center}

The sign of $w_{\mathrm{oa}}-w_{\mathrm{pc}}$ therefore decides whether the
outside air is a latent load or a latent resource. For the setter setpoint
$(\SI{37}{\celsius},\ \RH=0.50)$ the room humidity ratio is
$w_{\mathrm{pc}}=\SI{19.86}{\gram\per\kilogram}$, whose vapour pressure
corresponds to a dew point of \SI{24.9}{\celsius}: outdoor air is a
dehumidification load exactly when its dew point exceeds that. At the design
ambient (\SI{31}{\celsius}, $\RH=0.75$) it does ---
$w_{\mathrm{oa}}=\SI{21.36}{\gram\per\kilogram}$ --- but only by
\SI{1.5}{\gram\per\kilogram}. Gulf summer dew points straddle the criterion,
which is why the computed latent demand is small and sign-changeable rather than
overwhelming.

\subsection{The assembled demand, and where the ventilation went}
\label{sec:dem:assembled}

Collecting Eqs.~\eqref{eq:dem:Qeggs}--\eqref{eq:dem:mwext} gives the hatchery
generator exactly as coded:
\begin{align}
\qs&=N_{\mathrm{eggs}}q_{\mathrm{egg}}
+H_{\mathrm{tr}}\left(T_{\mathrm{sol}}-T_{\mathrm{pc}}\right)
+\dot m_{\mathrm{inf}}c_{p,\mathrm{a}}\left(T_{\mathrm{oa}}-T_{\mathrm{pc}}\right),
\label{eq:dem:qs_hatch}\\
\mdw&=N_{\mathrm{eggs}}r_{\mathrm{evap}}
+\dot m_{\mathrm{inf}}\left(w_{\mathrm{oa}}-w_{\mathrm{pc}}\right).
\label{eq:dem:mw_hatch}
\end{align}
(The code computes the latent term as a power using \verb|LOAD_HFG| and the
generator immediately divides it back out, so Eq.~\eqref{eq:dem:mw_hatch} is what
the interface actually receives --- an instance of the round trip discussed in
Section~\ref{sec:dem:hfg}, here exact because both conversions use the same
constant.)

\paragraph{The accounting point that must be right.} The mandatory fresh-air
ventilation does \emph{not} appear in
Eqs.~\eqref{eq:dem:qs_hatch}--\eqref{eq:dem:mw_hatch}, and this is not an
omission. In this architecture the fresh air enters the \emph{process loop} at
damper $d_{p\mathrm{A}}$ and is conditioned by the wheel and the cooler
\emph{before} it reaches the space. Its load is borne by the equipment and enters
the plant through the process-inlet mixing state, not through the space balance.
Counting it here as well would be double counting. Equations
\eqref{eq:dem:qs_hatch}--\eqref{eq:dem:mw_hatch} are the loads on the \emph{space
alone}: eggs, envelope, infiltration.

\paragraph{Where it does appear.} Combining the process-inlet mixing with the
required supply humidity, the moisture the wheel must actually remove is
\begin{equation}
\dot m_{\mathrm{H_2O}}^{\mathrm{wheel}}
=\mdp\left(w_{\mathrm{pi}}^{\mathrm{mix}}-w_{\mathrm{pi}}^{\mathrm{req}}\right)
=\underbrace{\mdv\left(w_{\mathrm{oa}}-w_{\mathrm{pc}}\right)}_{\text{forced ventilation}}
+\underbrace{\dot m_{\mathrm{evap}}
+\dot m_{\mathrm{inf}}\left(w_{\mathrm{oa}}-w_{\mathrm{pc}}\right)}_{\text{space sources}},
\label{eq:dem:trueduty}
\end{equation}
using the identity $\mdp\,d_{p\mathrm{A}}=\mdv$ guaranteed by the sizing of
Section~\ref{sec:dem:causality}. At the design point the three terms are
\SI{15.0}{\gram\per\second}, \SI{4.6}{\gram\per\second} and
\SI{4.4}{\gram\per\second}: a total wheel duty of
\SI{24.0}{\gram\per\second} against a \emph{space} latent load of only
\SI{9.0}{\gram\per\second}. \textbf{Nearly two thirds of the dehumidification
duty is imported by the ventilation air the eggs require.} The plant is sized by
a flow that cannot be reduced multiplied by a humidity gap set by the weather.

\subsection{Sizing: everything follows from \texorpdfstring{$N_{\mathrm{eggs}}$}{N\_eggs}}
\label{sec:dem:sizing}

Prescribing loads, flows and geometry independently allows them to describe
different-sized buildings at the same time --- an inconsistency that occurred
during development. The cure is a single declared capacity from which everything
descends.

\paragraph{Internal loads, linear in capacity.}
Equations~\eqref{eq:dem:Qeggs} and \eqref{eq:dem:mevap}.

\paragraph{Geometry, from packing density and a fixed ceiling height.}
\begin{equation}
A_{\mathrm{floor}}=\frac{N_{\mathrm{eggs}}}{n_{\mathrm{egg/m^2}}},\qquad
V_{\mathrm{room}}=A_{\mathrm{floor}}H_{\mathrm{room}},\qquad
A_{\mathrm{walls}}=4\sqrt{A_{\mathrm{floor}}}\,H_{\mathrm{room}},
\label{eq:dem:geom}
\end{equation}
the last assuming a square footprint, so the perimeter is
$4\sqrt{A_{\mathrm{floor}}}$.

\paragraph{Envelope, infiltration and forced ventilation.}
\begin{equation}
H_{\mathrm{tr}}=U_{\mathrm{env}}\left(2A_{\mathrm{floor}}+A_{\mathrm{walls}}\right),
\qquad
\dot m_{\mathrm{inf}}=\frac{\rho_{\mathrm{da}}V_{\mathrm{room}}\,\mathrm{ACH}_{\mathrm{inf}}}{3600},
\qquad
\mdv=\frac{\rho_{\mathrm{da}}N_{\mathrm{eggs}}v_{\mathrm{egg}}}{3600},
\label{eq:dem:Htr}
\end{equation}
the factor $2$ counting roof and floor.

\paragraph{Scaling laws.} How does the plant scale with capacity? Not as the
naive ``surface $\propto$ volume$^{2/3}$'', because \textbf{the ceiling height is
fixed}: the building spreads sideways rather than growing isotropically. From
Eqs.~\eqref{eq:dem:geom}--\eqref{eq:dem:Htr}, everything is linear in
$N_{\mathrm{eggs}}$ except the walls, which go as $\sqrt{N_{\mathrm{eggs}}}$:
\begin{equation}
H_{\mathrm{tr}}
=\underbrace{\frac{2U_{\mathrm{env}}}{n_{\mathrm{egg/m^2}}}N_{\mathrm{eggs}}}
_{\text{roof}+\text{floor},\ \propto N_{\mathrm{eggs}}}
+\underbrace{4U_{\mathrm{env}}H_{\mathrm{room}}
\sqrt{\frac{N_{\mathrm{eggs}}}{n_{\mathrm{egg/m^2}}}}}
_{\text{walls},\ \propto\sqrt{N_{\mathrm{eggs}}}}\;.
\label{eq:dem:scaling}
\end{equation}
Numerically, from $\num{150000}$ to $10^{6}$ eggs (a factor
$6.67$), $H_{\mathrm{tr}}$ rises from \SI{942}{\watt\per\kelvin} to
\SI{5495}{\watt\per\kelvin}, a factor $5.84$: an effective exponent of
$N_{\mathrm{eggs}}^{0.93}$, nearly but not exactly linear.

\begin{center}
\fbox{\parbox{0.92\textwidth}{
\textbf{Consequence: the load \emph{balance} is scale-invariant.} Since
essentially every term scales with $N_{\mathrm{eggs}}$, so do $\qs$, $\mdw$ and
therefore $\mdp$ --- and the \emph{ratios} (sensible-to-latent, and
$d_{p\mathrm{A}}=\mdv/\mdp$) do not move. Making the hatchery bigger does not
change whether it is latent- or sensible-dominated; that is set by the setpoint
and the climate. What does \emph{not} scale is the \textbf{well}: one borehole
delivers what it delivers, so raising $N_{\mathrm{eggs}}$ drives the plant into
the coil's large-flow limit. The binding constraint is the well-to-demand ratio,
not the building.}}
\end{center}

\subsection{Airflow causality --- which way the arrow points}
\label{sec:dem:causality}

\paragraph{The wrong way round.} An earlier version prescribed the fresh-air
damper and derived the supply flow from it, $\mdp=\mdv/d_{p\mathrm{A}}$. That
makes the fan flow a $1/d_{p\mathrm{A}}$ hyperbola: at $d_{p\mathrm{A}}=0.1$ it
returns \SI{100}{\kilogram\per\second} of supply air. This is an artifact of
inverted causality, not a physical result. A damper position is not a sizing
input.

\paragraph{The right way round.} The physical logic runs in four steps.
\begin{enumerate}[nosep]
\item The \textbf{ventilation is forced by the application}: embryos consume
oxygen and produce carbon dioxide, so $\mdv$ (Eq.~\eqref{eq:dem:Htr}) is the one
flow in the plant that is not a design choice.
\item The \textbf{supply flow follows the design sensible load} carried at a
prescribed depression --- the variable-volume closure of
Section~\ref{sec:dem:balances}, applied once at design time:
\begin{equation}
\mdp^{\mathrm{load}}=\frac{\qs^{\mathrm{design}}}
{c_{p,\mathrm{a}}\,\Delta T_{\mathrm{sup}}^{\mathrm{design}}},
\qquad \Delta T_{\mathrm{sup}}^{\mathrm{design}}=\SI{3}{\kelvin}.
\label{eq:dem:mdotload}
\end{equation}
\item The \textbf{ventilation is a floor}, and the damper is \emph{derived}:
\begin{equation}
\boxed{\;\mdp=\max\!\left(\mdp^{\mathrm{load}},\,\mdv\right),\qquad
d_{p\mathrm{A}}=\min\!\left(\frac{\mdv}{\mdp},\,1\right).\;}
\label{eq:dem:causality}
\end{equation}
The damper opens exactly far enough to admit the required fresh air and no
further.
\item If $\mdp^{\mathrm{load}}\le\mdv$ the ventilation dominates:
$d_{p\mathrm{A}}\to1$, recirculation vanishes and the plant runs
\SI{100}{\percent} once-through. The model raises a \verb|vent_dominated| flag;
this is the same condition the interface reports as \verb|vent_limited|.
\end{enumerate}

\paragraph{On the depression.} $\Delta T_{\mathrm{sup}}$ is the standard HVAC
sizing handle: small depression means a large fan and gentle, uniform conditions;
large depression means a small fan and a cold jet. \textbf{For a setter it must
be small} --- eggs cannot be chilled unevenly --- so hatcheries run high
recirculation at small temperature difference. The remaining hardware then
follows the streams it serves:
$\mdr=0.70\,\mdp$, $\md_{\mathrm{cool}}=1.5\,\mdp$ and
$\UA_{\mathrm{coil}}=\SI{3000}{\watt\per\kelvin}$ per \si{\kilogram\per\second}
of $\mdr$. Freezing those as constants while scaling the plant would be the same
class of error as the hyperbola above.

\subsection{The chain evaluated at the baseline}
\label{sec:dem:worked}

Table~\ref{tab:dem:chain} runs the whole descent at the notebook's baseline:
$N_{\mathrm{eggs}}=10^{6}$, room $(\SI{37}{\celsius},\ \RH=0.50)$, ambient
$(\SI{31}{\celsius},\ \RH=0.75)$.

\begin{table}[htbp]
\centering
\caption{The sizing chain and the interface solve at the baseline configuration.
Every input is read from the configuration cell; every output is the value the
notebook prints.}
\label{tab:dem:chain}
\begin{tabular}{@{}llr@{}}
\toprule
Step & Expression & Value \\
\midrule
\multicolumn{3}{@{}l}{\emph{Geometry and sizing}}\\
Egg sensible heat & $10^{6}\times\SI{0.13}{\watt}$ & \SI{130.0}{\kilo\watt}\\
Egg moisture      & $10^{6}\times\SI{4.6e-9}{\kilogram\per\second}$ & \SI{4.60}{\gram\per\second}\\
Floor area        & $10^{6}/200$ & \SI{5000}{\metre\squared}\\
Room volume       & $5000\times3.5$ & \SI{17500}{\metre\cubed}\\
Wall area         & $4\sqrt{5000}\times3.5$ & \SI{989.9}{\metre\squared}\\
Transmission      & $0.5\,(2\times5000+989.9)$ & \SI{5495}{\watt\per\kelvin}\\
Infiltration      & $1.2\times17500\times0.5/3600$ & \SI{2.917}{\kilogram\per\second}\\
Forced ventilation& $1.2\times(10^{6}\times0.030)/3600$ & \SI{10.00}{\kilogram\per\second}\\
\midrule
\multicolumn{3}{@{}l}{\emph{Design sensible load}, Eq.~\eqref{eq:dem:qs_hatch}, $T_{\mathrm{sol}}=\SI{35}{\celsius}$}\\
\quad eggs         &                                & $+\SI{130.0}{\kilo\watt}$\\
\quad envelope     & $5495\times(35-37)$            & $-\SI{11.0}{\kilo\watt}$\\
\quad infiltration & $2.917\times1006\times(31-37)$ & $-\SI{17.6}{\kilo\watt}$\\
\quad \textbf{total} &                              & $\mathbf{+\SI{101.4}{\kilo\watt}}$\\
\midrule
\multicolumn{3}{@{}l}{\emph{Design latent load}, Eq.~\eqref{eq:dem:mw_hatch}, $w_{\mathrm{oa}}-w_{\mathrm{pc}}=+\SI{1.50}{\gram\per\kilogram}$}\\
\quad eggs         &                                & $+\SI{4.60}{\gram\per\second}$\\
\quad infiltration & $2.917\times1.50\times10^{-3}$ & $+\SI{4.38}{\gram\per\second}$\\
\quad \textbf{total} &                              & $\mathbf{+\SI{8.98}{\gram\per\second}}$ ($=\SI{22.5}{\kilo\watt}$)\\
\midrule
\multicolumn{3}{@{}l}{\emph{Airflow}, Eqs.~\eqref{eq:dem:mdotload}--\eqref{eq:dem:causality}}\\
Load-driven flow  & $101405/(1006\times3.0)$ & \SI{33.60}{\kilogram\per\second}\\
$\mdp$            & $\max(33.60,\,10.00)$    & \SI{33.60}{\kilogram\per\second} (load wins)\\
$d_{p\mathrm{A}}$ & $10.00/33.60$            & $0.298$ (derived)\\
$\mdr$, $\md_{\mathrm{cool}}$ & $0.70\,\mdp$, $1.5\,\mdp$ & $23.52$, \SI{50.40}{\kilogram\per\second}\\
$\UA_{\mathrm{coil}}$ & $3000\times23.52$    & \SI{70.6}{\kilo\watt\per\kelvin}\\
\midrule
\multicolumn{3}{@{}l}{\emph{Interface solve}, constant-volume closure $\md_{\mathrm{sup}}=\mdp$, $c_p=\SI{1042.9}{\joule\per\kilogram\per\kelvin}$}\\
$T_{\mathrm{pi}}$ & $37-101405/(33.60\times1042.9)$ & \SI{34.11}{\celsius}\\
$w_{\mathrm{pi}}$ & $19.86-8.98/33.60$              & \SI{19.59}{\gram\per\kilogram}\\
$\Delta T_{\mathrm{sup}}$, $\Delta w_{\mathrm{sup}}$ & & \SI{2.89}{\kelvin}, \SI{0.267}{\gram\per\kilogram}\\
$\RH_{\mathrm{pi}}$, SHR & & $0.58$, $0.82$\\
\bottomrule
\end{tabular}
\end{table}

\paragraph{Three things to read off the table.} First, \textbf{the eggs are more
than the whole load}: at \SI{31}{\celsius} ambient the room at \SI{37}{\celsius}
is warmer than outside, so the envelope and the infiltration are both
\emph{negative} --- they help --- and the \SI{130}{\kilo\watt} of metabolic heat
is $128\,\%$ of the \SI{101.4}{\kilo\watt} net. The building is an internally
driven load in a mild-by-comparison climate, and $q_{\mathrm{egg}}$ is by a wide
margin the highest-leverage placeholder in the model.

Second, \textbf{the achieved depression is not the design depression}:
\SI{2.89}{\kelvin} against \SI{3.0}{\kelvin}. The flow is sized with the dry-air
$c_{p,\mathrm{a}}=\SI{1006}{\joule\per\kilogram\per\kelvin}$ while the interface
solves with the moist-air value $\SI{1042.9}{\joule\per\kilogram\per\kelvin}$ at
$w_{\mathrm{pc}}$. The $3.7\,\%$ discrepancy is harmless but real, and it is
exactly the kind of double-definition the moisture-rate argument of
Section~\ref{sec:dem:hfg} was written to avoid.

Third, \textbf{the required conditioning is tiny in humidity terms}:
\SI{0.267}{\gram\per\kilogram} of depression. The room condition line is very
steep and the plant's difficulty lies almost entirely in the ventilation air it
must dry on the way in (Eq.~\eqref{eq:dem:trueduty}), not in the space itself.

\subsection{The interface contract: what a second demand case must supply}
\label{sec:dem:contract}

The point of the preceding section is that none of it was privileged. To attach
any other application to this plant, the following --- and only the
following --- must be provided.

\begin{enumerate}[nosep]
\item \textbf{A load generator.} A function of the ambient state (and of
whatever internal state the application has) returning the pair
$(\qs,\mdw)$: sensible load in \si{\watt}, moisture load in
\si{\kilogram\per\second}, both positive when the plant must remove. If the
application only knows kilowatts, it must convert once, explicitly, at its own
boundary.
\item \textbf{A room setpoint} $(T_{\mathrm{pc}},w_{\mathrm{pc}})$, with
$w_{\mathrm{pc}}$ obtained from the setpoint relative humidity through the
psychrometric relations, because both balances are written about the room state.
\item \textbf{Exactly one closure}: a supply flow (constant volume) or a supply
depression (variable volume), optionally with a minimum flow $\md_{\min}$
representing a mandatory ventilation rate. Not both, not neither.
\item \textbf{No double counting.} Any stream the plant conditions upstream of
the space --- ventilation air admitted at $d_{p\mathrm{A}}$ being the case in
point --- must be excluded from the generator. The generator returns loads on the
\emph{space}, not on the plant.
\item \textbf{Attention to the flags.} \verb|vent_limited|, and the two guards of
Section~\ref{sec:dem:guards}, are warnings by design; a caller that ignores them
will receive a supply state that no equipment can produce and will not be told
twice.
\end{enumerate}

\paragraph{One honest caveat.} The interface is generic at \emph{run} time, but
the plant's \emph{sizing} constants --- $\mdp$, $\mdr$, $\md_{\mathrm{cool}}$,
$\UA_{\mathrm{coil}}$ --- still descend from the hatchery chain of
Section~\ref{sec:dem:sizing}. A second application inherits the interface for
free; to inherit the hardware it must also declare its own design sensible load
and design depression and re-run Eqs.~\eqref{eq:dem:mdotload}--\eqref{eq:dem:causality}.
Generalising the demand side was the first half of that job; generalising the
sizing chain is the second, and it is not yet done.
